% =============================================================
%  ch15_projekt1.tex  --  Projekt 1: Stepper-Roboter mit Trajektorienbefehl
% =============================================================

\chapter{Projekt 1 -- Stepper-Roboter mit Trajektorienbefehl vom PC}

\textbf{Ziel:} Vom externen PC eine gewünschte Trajektorie senden, der Roboter
führt sie aus. Da es ein langer Auftrag mit Fortschritt/Abbruch ist, ist die
natürliche Schnittstelle eine \textbf{Action}.

\section{ROS\,2-Bausteine}
\begin{longtable}{@{}p{2.6cm}p{12cm}@{}}
\toprule
\textbf{Typ} & \textbf{Konkret} \\
\midrule
Nodes   & \code{trajectory\_commander} (PC, Action-Client) $\cdot$
          \code{controller\_manager} + \code{joint\_trajectory\_controller} +
          \code{joint\_state\_broadcaster} $\cdot$
          Hardware-Interface (Sim: \code{gz\_ros2\_control};
          Real: micro-ROS$\rightarrow$ESP32) \\
Actions & \code{control\_msgs/FollowJointTrajectory}
          (Goal: Bahn; Feedback: Ist-Position; Result: Erfolg) \\
Topics  & \code{/joint\_states} (Encoder/StallGuard) $\cdot$
          \code{/joint\_trajectory\_controller/\ldots} $\cdot$ \code{/tf} \\
Services & Controller laden/umschalten
           (\code{/controller\_manager/switch\_controller}) \\
Parameter & Gelenklimits, Getriebeübersetzung, Reglerverstärkungen \\
\bottomrule
\end{longtable}

\section{rqt\_graph (vereinfacht)}
\begin{center}
\begin{tikzpicture}[node distance=6mm and 16mm]
\node[node] (cmd)  {/trajectory\_commander\\\scriptsize (PC, Action-Client)};
\node[node,right=of cmd] (jtc)  {/joint\_trajectory\_\\controller};
\node[node,right=of jtc] (hw)   {/hardware\_interface\\\scriptsize (ESP32 via micro-ROS)};
\node[topic,below=10mm of jtc] (js) {/joint\_states};
\node[node,below=6mm of cmd]   (rviz) {/rviz2};
\draw[flow] (cmd) -- node[topic,fill=white,inner sep=1pt]{\scriptsize action goal} (jtc);
\draw[flow] (jtc) -- node[topic,fill=white,inner sep=1pt]{\scriptsize commands} (hw);
\draw[flow] (hw)  |- (js);
\draw[flow] (js)  -| (jtc);
\draw[flow] (js)  -| (rviz);
\draw[flow,dashed,accent2] (jtc) to[bend left=20]
  node[above,font=\scriptsize\itshape]{feedback} (cmd);
\end{tikzpicture}
\captionof{figure}{Projekt~1: Der PC schickt eine Trajektorie als Action-Goal;
der Controller kommandiert das Hardware-Interface (ESP32) und liest
\code{/joint\_states} zur Regelung zurück.}
\end{center}

\noindent\textbf{Datenfluss:}
PC $\rightarrow$ Action-Goal (Bahn) $\rightarrow$
\code{joint\_trajectory\_controller} interpoliert $\rightarrow$
Sollwerte ans Hardware-Interface $\rightarrow$
ESP32 erzeugt Schritte am TMC2209, liest StallGuard/Encoder $\rightarrow$
\code{/joint\_states} zurück $\rightarrow$ Regelschleife schließt sich.
In Simulation ersetzt \code{gz\_ros2\_control} den ESP32 -- der PC-Code bleibt
unverändert.
